\documentclass[a4paper]{article}
\usepackage[pdftex]{graphicx}
\usepackage[utf8]{inputenc}
\usepackage{enumerate}
\usepackage{icomma}
\usepackage{href-ul}
\hypersetup{
	colorlinks=true,
	linkcolor=blue,
	urlcolor=blue}
\usepackage{siunitx}
\sisetup{locale=DE}
\usepackage{amssymb}
\usepackage{geometry}
\geometry{a4paper, top=15mm, left=15mm, right=15mm, bottom=15mm,
	headsep=10mm, footskip=12mm}

\begin{document}
	
	\begin{enumerate}[1.]
		
		%Task 1
		{\bf \item When Galileo Galilei wanted to determine the speed of light, it was not successful.}
		\begin{enumerate}[a)]
			\item Describe in bullet points a measuring method that can be used to determine the speed of light.
			\item Specify the symbol and value of the speed of light.
		\end{enumerate}
		
		%Exercise 2
		{\bf \item The light from two flashlights $L_1(1|2)$ and $L_2(2|6)$ hits a body,} which starts at point $A(6,5|2)$ and at point $ B(6.5|4)$ ends. Behind the body there is a wall that forms the line with the points $S_1(8|0)$ and $S_2(8|6)$.
		
		\begin{enumerate}[a)]
			\item Draw the points in a coordinate system and construct the shadow of the body
			\item Divide the shadow into different areas and use technical terms!
		\end{enumerate}
		
		%Task 3
		{\bf \item Specify the name of the lens shown.}
		
		\includegraphics[width=7 cm]{linsen100.pdf}
		
		%Task 4
		{\bf \item Kepler's and the Dutch (Galilean) telescope are probably among the best known. }
		List three key differences between these two.
		
		%Task 5
		{\bf \item A converging lens with the focal length $f=\SI{40}{\milli\meter}$ is available.}
		
		\begin{enumerate}[a)]
			\item Construct with the object width $g=\SI{15}{\milli\meter}$ and the object size
			$G=\SI{20}{\milli\meter}$ the image size $B$.
			\item Indicate the image that results from this and name two properties of this image.
		\end{enumerate}
		
		%Task 6
		{\bf \item Explain how myopia occurs and how it can be corrected.}
		
		%Task 7
		{\bf \item Two of the invisible components of sunlight are ultraviolet light and infrared light.} Give three examples of how ultraviolet radiation can be used.
		
		%Task 8
		{\bf \item Tick the relevant statements about dispersion:}
		
		$\square$ The fanning out of the spectral colors is called dispersion.\\
		$\square$ Monochromatic light always produces a line spectrum.\\
		$\square$ UV light and infrared light are not spectral colors.\\
		$\square$ White sunlight is an example of monochromatic light.\\
		$\square$ A continuous spread of light is called a continuous spectrum.
		
		%Task 9
		{\bf \item task}\\
		\begin{minipage}{0.5\textwidth}
			\begin{enumerate}[a)]
				\item Explain in brief what is meant by total internal reflection.
				\item Draw the beam path through the prism.
			\end{enumerate}
		\end{minipage}
		\hfill
		\begin{minipage}{0.45\textwidth}
			\includegraphics[width=5 cm]{prisma100.pdf}
		\end{minipage}
		
		%Task 10
		{\bf \item A firmly clamped converging lens creates a real image of a luminous object on a frosted glass screen.} What do you have to change if you want to get a smaller, sharp image of the object?
	\end{enumerate}
	
	\fbox{
		\begin{minipage}{0.5\textwidth}
			Additional worksheets are available at
			
			\href{http://www.okuyakl.com}{www.okuyakl.com}
		\end{minipage}
		\hfill
		\begin{minipage}{0.4\textwidth}
			\includegraphics[width=1.5 cm]{../../viecher/zwe03}
			\includegraphics[width=2 cm]{../../viecher/afanticon1}
			
	\end{minipage}}
	
\end{document}